\appendix\label{appendix}
\addchap{Explanations of Abbreviations}

This chapter provides an index of morpheme glosses and their respective morpheme category as categorized for the Navajo Corpus of Historical Narratives. The glosses include abbreviated and non-abbreviated morphemes that are not part of the root. In addition, the different allomorphs found in the corpus are listed to facilitate the identification in the texts.

\section*{Features}\label{appendix:features}

This section explains the abbreviated morphological glosses for morphs in the morphological categories. Only general features have been abbreviated. Therefore, if a morph's semantics was too specific it has been written out explicitly, like \textit{j-} `into\_distance'. In other cases, if the semantics were too dependent on the meaning of the root or not identifiable, it has been decided to use the gloss \textsc{th} for ``thematic''. The entry indicates the meaning of the gloss, together with the category of the word root that they belong to (``v'' if verbal, ``pro'' if pronominal, ``n'' if nominal, and ``post'' if postpositional) and, additionally for verbs, their morphological position (v:\textsc{prvb} `preverb'; v:\textsc{iter} `iterative'; v:\textsc{distr} `distributive'; v:\textsc{obj} `object'; v:\textsc{ar} `areal'; v:\textsc{osbj} `outer subject'; v:\textsc{qual} `qualifier'; v:\textsc{ams} `aspect, mood, and subject'; v:\textsc{val} `valence', v `verbal root', see the entire template in \figref{fig:verbtemplate}. It should be noted that some morphs express several of these features simultaneously, such as 1.\textsc{sg}.\textsc{si}.\textsc{perf}.

\begin{xltabular}{\textwidth}{lXX}
\caption{Abbreviations and features} \label{tab:abbrevfeat}\\

    \lsptoprule
        \textsc{gloss} & \textsc{category} & \textsc{position}\\
    \midrule
    \endfirsthead
    \toprule
        \textsc{gloss} & \textsc{category} & \textsc{position}\\
	\midrule
    \endhead

    \endfoot  
    \endlastfoot
1.\textsc{du}/\textsc{pl} & first person duoplural & (v:\textsc{obj}; v:\textsc{ams}; v, attaches to any category, pro)\\
1.\textsc{sg} & first person singular & (v:\textsc{obj}; v:\textsc{ams}; v, attaches to any category)\\
2.\textsc{du}/\textsc{pl} & second person duoplural & (v:\textsc{obj}; v:\textsc{ams}; v, attaches to any category, pro)\\
2.\textsc{sg} & second person singular & (v:\textsc{obj}; v:\textsc{ams}; v, attaches to any category, pro)\\
3 & third person & (pro)\\
3.\textsc{bi} & third person \textit{bi-} & (v:\textsc{obj}; attaches to any category)\\
3.\textsc{yi} & third person \textit{yi-} & (v:\textsc{obj}; attaches to any category)\\
\textsc{abs} & absolute aspect & (v)\\
\textsc{agt}.\textsc{pas} & agentive passive & (v:\textsc{osbj})\\
\textsc{anim} & animate person & (v:\textsc{osbj}; v:\textsc{obj}, v, pro)\\
\textsc{ano} & animate object & (v)\\
\textsc{ar} & areal & (v:\textsc{ar})\\
\textsc{ar}.\textsc{ni}.\textsc{th} & areal with \textit{ni-}thematic contraction (from the \textsc{qual} position) & (v:\textsc{ar})\\
\textsc{comp} & comparative aspect & (v)\\
\textsc{con} & conative aspect & (v)\\
\textsc{conc} & conclusive aspect & (v)\\
\textsc{cont} & continuative aspect & (v)\\
\textsc{curs} & cursive aspect & (v)\\
\textsc{d} & d-classifier (includes \mbox{d-segment} of first person duoplural forms) & (v:\textsc{val})\\
\textsc{dir} & directive & (v:\textsc{qual})\\
\textsc{distr} & distributive & (v:\textsc{distr}, attaches to any category; v)\\
\textsc{distr2} & distributive 2 & (v)\\
\textsc{du} & dual & (v)\\
\textsc{du}/\textsc{pl} & dual or plural & (v)\\
\textsc{dur} & durative & (v)\\
\textsc{ffo} & flat flexible object & (v)\\
\textsc{foc} & focus & (prt)\\
\textsc{fut} & future & (v:\textsc{qual}, v, prt)\\
\textsc{gen} & generic & (prt)\\
\textsc{impf} & imperfective & (v)\\
\textsc{inc} & inceptive & (v:\textsc{qual})\\
\textsc{inch} & inchoative & (v:\textsc{qual})\\\textsc{indef} & indefinite & (prt)\\
\textsc{intrg} & interrogative & (prt)\\
\textsc{ints} & intensifier & (v)\\
\textsc{iter} & iterative & (v:\textsc{iter}, v)\\
\textsc{l} & l-classifier (includes voicing of ł-classifier) & (v:\textsc{val})\\
\textsc{ł} & ł-classifier (includes devoicing of l-classifier) & (v:\textsc{val})\\
\textsc{lpb} & load, pack, burden & (v)\\
\textsc{mom}1 & momentaneous aspect (1) & (v)\\
\textsc{mom}2 & momentaneous aspect (2) & (v)\\
\textsc{mom}3 & momentaneous aspect (3) & (v)\\
\textsc{neu} & neuter aspect & (v)\\
\textsc{neu}1 & neuter aspect (1) & (v)\\
\textsc{ni}.\textsc{impf} & ni-imperfective & (v:\textsc{ams})\\
\textsc{ni}.\textsc{perf} & ni-perfective & (v:\textsc{ams})\\
\textsc{nmlz} & nominalizer & (prt)\\
\textsc{nsp}.\textsc{obj} & non-specific object & (v:\textsc{obj})\\
\textsc{nsp}.\textsc{sbj} & non-specific subject & (v:\textsc{sbj})\\
\textsc{oc} & open container & (v)\\
\textsc{opt} & optative & (v:\textsc{ams}, v)\\
\textsc{pas} & passive & (v)\\
\textsc{pcl} & particular&\\
\textsc{pcl}.\textsc{ref} & particular reference\\
\textsc{perf} & perfective & (v)\\
\textsc{pl} & plural & (v)\\
\textsc{poss} & possessor & (attaches to any category)\\
\textsc{pres}.\textsc{ref} & present referential & (prt)\\
\textsc{prog} & progressive & (v:\textsc{ams}, v)\\
\textsc{pst} & past & (prt)\\
\textsc{pst}.\textsc{ref} & past referential & (prt)\\
\textsc{rec} & reciprocal & (v:\textsc{obj}; attaches to any category)\\
\textsc{refl} & reflexive & (v:\textsc{obj}; v:\textsc{ams}; v, attaches to any category)\\
\textsc{rep} & repetitive & (v:\textsc{prvb}, v)\\
\textsc{rev} & reversionary & (v:\textsc{prvb}, v)\\
\textsc{semf} & semelfactive & (v:\textsc{qual}, v)\\
\textsc{ser} & seriative & (v:\textsc{qual})\\
\textsc{sfo} & slender flexible object & (v)\\
\textsc{sg} & singular & (v)\\
\textsc{si}.\textsc{impf} & \textit{si-}imperfective & (v:\textsc{ams})\\
\textsc{si}.\textsc{perf} & \textit{si-}perfective & (v:\textsc{ams})\\
\textsc{sitr} & semeliterative & (v:\textsc{prvb})\\
\textsc{sp} & specific & (prt)\\
\textsc{sro} & solid roundish object & (v)\\
\textsc{sso} & stiff slender object & (v)\\
\textsc{sub} & subordination & (prt)\\
\textsc{term} & terminative & (v:\textsc{qual})\\
\textsc{th} & thematic (unclear lexical \mbox{semantics}) & (v:\textsc{prvb}; v:\textsc{qual})\\
\textsc{trnsl} & transitional & (v:\textsc{qual}, v)\\
\textsc{usit} & usitative aspect & (v)\\
\textsc{yi}.\textsc{perf} & \textit{yi-}perfective & (v:\textsc{ams})\\
\textsc{yi}.\textsc{zero}.\textsc{perf} & \textit{yi-}zero perfective & (v:\textsc{ams})\\
\textsc{zero}.\textsc{impf} & zero-imperfective & (v:\textsc{ams})\\
\lspbottomrule
\end{xltabular}

\section*{Verbal glosses, their categories and morphs}\label{appendix:verbglosses}

The following table provides a detailed list of the different verbal prefixes in the narratives, with reference to their categories and list of morphs. Given the high amount of homonymy, allomorphy and syncretism across morphemes, the reader might find this table helpful in identifying the glosses, especially for the ones labeled as \textsc{th} (thematic). These include also non-abbreviated verbal glosses in affixes, but excludes glosses for root morphemes.

\begin{xltabular}{\textwidth}{llX}
\caption{Verbal prefixes in the narratives} \label{tab:abbrevverb}\\

    \lsptoprule
        \textsc{gloss} & \textsc{position} & \textsc{morphs}\\
    \midrule
    \endfirsthead
    \toprule
        \textsc{gloss} & \textsc{position} & \textsc{morphs}\\
	\midrule
    \endhead

    \endfoot  
    \endlastfoot
 1.\textsc{du}/\textsc{pl} & \textsc{obj} & \textit{nihi} (\textit{nih}, \textit{nah})\\                 
 1.\textsc{du}/\textsc{pl}.\textsc{ni}.\textsc{impf} & \textsc{ams} &  \textit{nii }(\textit{níi})\\
 1.\textsc{du}/\textsc{pl}.\textsc{ni}.\textsc{perf} & \textsc{ams} &  \textit{nii} (\textit{nii})\\
 1.\textsc{du}/\textsc{pl}.\textsc{prog} & \textsc{ams} &  \textit{ii} (\textit{yii}, \textit{íi}, \textit{i})\\
 1.\textsc{du}/\textsc{pl}.\textsc{si}.\textsc{perf} & \textsc{ams} &  \textit{sii} (\textit{shii})\\
 1.\textsc{du}/\textsc{pl}.\textsc{yi}.\textsc{perf} & \textsc{ams} &  \textit{ii} (\textit{yii}, \textit{íi}, \textit{i})\\
 1.\textsc{du}/\textsc{pl}.\textsc{yi}.\textsc{zero}.\textsc{imp} & \textsc{ams} &  \textit{ii}\\
 1.\textsc{du}/\textsc{pl}.\textsc{yi}.\textsc{zero}.\textsc{perf} & \textsc{ams} &  \textit{ii} (\textit{ii})\\
 1.\textsc{du}/\textsc{pl}.\textsc{zero}.\textsc{impf} & \textsc{ams} &  \textit{ii} (\textit{i}, \textit{yii}, \textit{íi})\\
 1.\textsc{sg} & \textsc{obj} & shi (\textit{sh}, \textit{shí}, \textit{si})\\
 1.\textsc{sg}.\textsc{ni}.\textsc{impf} & \textsc{ams} &  \textit{nish} (\textit{nish})\\
 1.\textsc{sg}.\textsc{ni}.\textsc{perf} & \textsc{ams} &  \textit{ní} (\textit{nísh}, \textit{nís})\\
 1.\textsc{sg}.\textsc{prog} & \textsc{ams} &  \textit{é} (\textit{ées}, \textit{éesh}, \textit{yish})\\
 1.\textsc{sg}.\textsc{si}.\textsc{perf} & \textsc{ams} &  \textit{é} (\textit{sé})\\
 1.\textsc{sg}.\textsc{yi}.\textsc{perf} & \textsc{ams} &  \textit{ii} (\textit{áás}, \textit{íish}, \textit{aas}, \textit{í}, \textit{yí}, \textit{éé}, \textit{eesh})\\
 1.\textsc{sg}.\textsc{yi}.\textsc{zero}.\textsc{imp} & \textsc{ams} &  \textit{iis} (\textit{iish})\\
 1.\textsc{sg}.\textsc{yi}.\textsc{zero}.\textsc{perf} & \textsc{ams} &  \textit{ii} (\textit{oo}, \textit{iish}, \textit{yii}, \textit{i}, \textit{ash})\\
 1.\textsc{sg}.\textsc{zero}.\textsc{impf} & \textsc{ams} &  \textit{sh} (\textit{s}, \textit{yish})\\
 2.\textsc{du}/\textsc{pl} & \textsc{obj} & \textit{nihi}\\
 2.\textsc{du}/\textsc{pl}.\textsc{ni}.\textsc{impf} & \textsc{ams} & \textit{nó} (\textit{nó})\\
 2.\textsc{du}/\textsc{pl}.\textsc{prog} & \textsc{ams} & oo (\textit{oo})\\
 2.\textsc{du}/\textsc{pl}.\textsc{si}.\textsc{perf} & \textsc{ams} & \textit{soo} (\textit{sooh})\\
 2.\textsc{du}/\textsc{pl}.\textsc{yi}.\textsc{perf} & \textsc{ams} & \textit{aah} (\textit{aah})\\
 2.\textsc{du}/\textsc{pl}.\textsc{yi}.\textsc{zero}.\textsc{impf}& \textsc{ams} &\textit{ooh}\\
 2.\textsc{du}/\textsc{pl}.\textsc{zero}.\textsc{impf} & \textsc{ams} & \textit{oh} (\textit{o}, \textit{h}, \textit{ó}, \textit{óh})\\
 2.\textsc{sg} & \textsc{obj} & \textit{ni} (\textit{n})\\
 2.\textsc{sg}.\textsc{ni}.\textsc{impf} & \textsc{ams} & \textit{ní} (\textit{nín})\\
 2.\textsc{sg}.\textsc{opt} & \textsc{ams} & \textit{óó} (\textit{óó})\\
 2.\textsc{sg}.\textsc{prog} & \textsc{ams} & \textit{yí} (\textit{yí})\\
 2.\textsc{sg}.\textsc{si}.\textsc{perf} & \textsc{ams} & \textit{síní} (\textit{shíní})\\
 2.\textsc{sg}.\textsc{yi}.\textsc{perf} & \textsc{ams} & \textit{ini} (\textit{ini})\\
 2.\textsc{sg}.\textsc{yi}.\textsc{zero}.\textsc{impf} & \textsc{ams} &\textit{iin}\\
 2.\textsc{sg}.\textsc{zero}.\textsc{impf} & \textsc{ams} & \textit{ni} (\textit{í}, \textit{ó}, \textit{n}, \textit{ń}, \textit{i}, \textit{ín})\\
 3.\textsc{bi} & \textsc{obj} & \textit{bi} (\textit{b}, \textit{bí})\\
 3.\textsc{ni}.\textsc{impf} & \textsc{ams} & \textit{ni} (\textit{ee}, \textit{í}, \textit{yí}, \textit{ée}, \textit{é}, \textit{é})\\
 3.\textsc{ni}.\textsc{perf} & \textsc{ams} & \textit{ni} (\textit{ní}, \textit{ee}, \textit{í}, \textit{yí}, \textit{ń}, \textit{á}, \textit{é}, \textit{oo})\\
 3.\textsc{opt} & \textsc{ams} & \textit{óo} (\textit{óo})\\
 3.\textsc{opt} & \textsc{ams} & \textit{wo} (\textit{wo})\\
 3.\textsc{prog} & \textsc{ams} & \textit{oo} (\textit{yoo}, \textit{yíí}, \textit{oo}, \textit{oo}, \textit{oo}, \textit{oo}, \textit{oo}, \textit{oo}, \textit{oo}, \textit{oo})\\
 3.\textsc{si}.\textsc{impf} & \textsc{ams} & \textit{z} (\textit{z})\\
 3.\textsc{si}.\textsc{perf} & \textsc{ams} & \textit{si} (\textit{sh}, \textit{zh}, \textit{z}, \textit{s}, \textit{eez}, \textit{eesh}, \textit{ees}, \textit{éé}, \textit{shi}, \textit{yish}, \textit{éés}, \textit{ées}, \textit{éesh}, \textit{yis}, \textit{eezh}, \textit{h}, \textit{ee}, \textit{áaz}, \textit{íí})\\
 3.\textsc{yi} & \textsc{obj} & \textit{yi} (\textit{í}, \textit{y}, \textit{i})\\
 3.\textsc{yi}.\textsc{perf} & \textsc{ams} & \textit{yi} (\textit{oo}, \textit{yíí}, \textit{yoo}, \textit{yi}, \textit{yóó}, \textit{yéé}, \textit{yó}, \textit{yaa}, \textit{yaas}, \textit{yiz}, \textit{yii}, \textit{ii})\\
 3.\textsc{yi}.\textsc{zero}.\textsc{impf} & \textsc{ams} &\textit{ii} (\textit{i}, \textit{yii}, \textit{oo})\\
 3.\textsc{yi}.\textsc{zero}.\textsc{perf} & \textsc{ams} & \textit{yi} (\textit{yoo}, \textit{yii}, \textit{íi}, \textit{aa}, \textit{a}, \textit{i})\\
 3.\textsc{zero}.\textsc{impf} & \textsc{ams} & \textit{yi}\\
 against\_around & \textsc{prvb} & \textit{éé’}\\                   
 \textsc{agt}.\textsc{pas} & \textsc{osbj} & \textit{’di} (\textit{’d})\\
 \textsc{anim} & \textsc{obj} & \textit{ha} (\textit{hw}, \textit{h}, \textit{ho}, \textit{há})\\
 \textsc{anim} & \textsc{osbj} & \textit{ji} (\textit{zh}, \textit{j}, \textit{dzi}, \textit{sh}, \textit{z}, \textit{dz}, \textit{jí}, \textit{dzí})\\
 \textsc{ar} & \textsc{ar} & \textit{ha} (\textit{h}, \textit{ho}, \textit{hw}, \textit{há})\\
 \textsc{ar}.\textsc{ni}.\textsc{th} & \textsc{ar} & \textit{hó} (\textit{hoo})\\
 away\_from & \textsc{prvb} & \textit{gha}\\                       
 away\_from\_water & \textsc{prvb} & \textit{tsíts’á}\\                   
 block & \textsc{prvb} & \textit{dáá} (\textit{dá})\\
 \textsc{cont} & \textsc{prvb} & \textit{na} (\textit{naa}, \textit{n}, \textit{ni}, \textit{ne})\\
 \textsc{d} & \textsc{val} & \textit{d} (\textit{t,  ’})\\
 \textsc{dir} & \textsc{qual} & \textit{ííní-} (\textit{íí}, \textit{yó}, \textit{ii}, \textit{yíní}, \textit{íín}, \textit{ín}, \textit{i}, \textit{oo}, \textit{woo}, \textit{yin}, \textit{y}, \textit{wo}, \textit{óo}, \textit{ó}, \textit{wó}, \textit{o}, \textit{óó})\\
 \textsc{distr} & \textsc{distr} & \textit{da} (\textit{daa}, \textit{de})\\
 downward & \textsc{prvb} & \textit{’ada}\\
 downward & \textsc{prvb} & \textit{na} (\textit{n})\\
 friend & \textsc{prvb} & \textit{ch’oo}\\
 \textsc{fut} & \textsc{qual} & \textit{d} (\textit{dí}, \textit{di})\\
 half & \textsc{prvb} & \textit{’ałní} (\textit{’ałné})\\
 in\_a.circle & \textsc{prvb} & \textit{’ahéé}\\
 in\_two & \textsc{prvb} & \textit{k’í}\\
 \textsc{inc} & \textsc{qual} & \textit{di} (\textit{d}, \textit{dí})\\
 \textsc{inch} & \textsc{qual} & \textit{’n}\\
 into\_distance & \textsc{prvb} & \textit{j}\\
 \textsc{iter} & \textsc{iter} & \textit{ná} (\textit{né}, \textit{ń}, \textit{n}, \textit{náá}, \textit{ní}, \textit{é})\\
 \textsc{l} & \textsc{val} & \textit{l}\\
 \textsc{ł} & \textsc{val} & \textit{ł} (\textit{l})\\
 meaning & \textsc{prvb} & \textit{’ááł}\\
 \textsc{nsp} & \textsc{obj} & \textit{’a} (\textit{’}, \textit{’í})\\
 \textsc{nsp}.\textsc{sbj} & \textsc{osbj} & \textit{’á} (\textit{’á}, \textit{’}, \textit{’í})\\
 obtain & \textsc{prvb} & \textit{shó}\\
 out & \textsc{prvb} & \textit{ch’í} (\textit{ch’}, \textit{ch’é})\\
 over & \textsc{prvb} & \textit{naa}\\
 peace & \textsc{prvb} & \textit{k’é}\\
 plant & \textsc{prvb} & \textit{k’é}\\
 ready & \textsc{prvb} & \textit{hasht’ee} (\textit{hasht’e})\\
 \textsc{rec} & \textsc{obj} & \textit{’ahi} (\textit{’ah}, \textit{áhi}, \textit{ahi}, \textit{hi}, \textit{hi’})\\
 \textsc{refl} & \textsc{obj} & \textit{’ádi} (\textit{di})\\
 \textsc{rep} & \textsc{prvb} & \textit{náá} (\textit{ná}, \textit{éé}, \textit{ní}, \textit{né}, \textit{á}, \textit{áá})\\
 \textsc{rev} & \textsc{prvb} & \textit{ná} (\textit{ń}, \textit{ní}, \textit{náá}, \textit{áá}, \textit{éé}, \textit{né}, \textit{á})\\
 \textsc{semf} & \textsc{qual} & \textit{i} (\textit{oo}, \textit{yi})\\
 \textsc{ser} & \textsc{qual} & \textit{h} (\textit{ii}, \textit{hi}, \textit{ha}, \textit{há}, \textit{y}, \textit{haa}, \textit{yii}, \textit{i}, \textit{hé})\\
 \textsc{sitr} & \textsc{prvb} & \textit{náá} (\textit{náán}, \textit{nááná}, \textit{nááné}, \textit{nán}, \textit{éé})\\
 so & \textsc{prvb} & \textit{’ákó} (\textit{’áko}, \textit{kó}, \textit{’ákwí}, \textit{’ákwíi}, \textit{’akó}, \textit{’áho})\\
 so & \textsc{prvb} & \textit{’á}\\
 start & \textsc{prvb} & \textit{’adi}\\
 start & \textsc{prvb} & \textit{niki} (\textit{nik})\\
 stop & \textsc{prvb} & \textit{nii}\\
 straight\_out & \textsc{prvb} & \textit{k’i}\\
 suffering & \textsc{prvb} & \textit{’atí} (\textit{ti’})\\
 suffering & \textsc{prvb} & \textit{ti’}\\
 \textsc{term} & \textsc{qual} & \textit{ni} (\textit{n}, \textit{y})\\
 \textsc{th} & \textsc{obj} & \textit{’a} (\textit{’a}, \textit{’}, \textit{’i})\\
 \textsc{th} & \textsc{qual} & \textit{di} (\textit{d}, \textit{dí})\\
 \textsc{th} & \textsc{qual} & \textit{dí} (\textit{dée})\\
 \textsc{th} & \textsc{qual} & \textit{dz}\\
 \textsc{th} & \textsc{qual} & \textit{dzó}\\
 \textsc{th} & \textsc{qual} & \textit{h} (\textit{hé}, \textit{ha})\\
 \textsc{th} & \textsc{qual} & \textit{hí} (\textit{í}, \textit{íi}, \textit{háa})\\
 \textsc{th} & \textsc{qual} & \textit{hi} (\textit{y}, \textit{yii}, \textit{ii})\\
 \textsc{th} & \textsc{qual} & \textit{iiní}\\
 \textsc{th} & \textsc{qual} & \textit{j}\\
 \textsc{th} & \textsc{qual} & \textit{ł} (\textit{łi})\\
 \textsc{th} & \textsc{qual} & \textit{ni} (\textit{n},\textit{ í}, \textit{ní}, \textit{ííní})\\
 \textsc{th} & \textsc{qual} & \textit{ni} (\textit{oo}, \textit{ée}, \textit{n})\\
 \textsc{th} & \textsc{qual} & \textit{nó}\\
 \textsc{th} & \textsc{qual} & \textit{oo}\\
 \textsc{th} & \textsc{qual} & \textit{sh}\\
 \textsc{th} & \textsc{qual} & \textit{si} (\textit{ii}, \textit{s}, \textit{z}, \textit{y})\\
 \textsc{th} & \textsc{qual} & \textit{so}\\
 \textsc{th} & \textsc{qual} & \textit{y} (\textit{ii}, \textit{i}, \textit{oo}, \textit{yí}, \textit{yii}, \textit{yi})\\
 \textsc{th} & \textsc{prvb} & \textit{dá}\\
 \textsc{th} & \textsc{prvb} & \textit{da} (\textit{daa})\\
 \textsc{th} & \textsc{prvb} & \textit{dzi}\\
 \textsc{th} & \textsc{prvb} & \textit{ha} (\textit{haa}, \textit{hw}, \textit{h}, \textit{há})\\
 \textsc{th} & \textsc{prvb} & \textit{kéé} (\textit{ké})\\
 \textsc{th} & \textsc{prvb} & \textit{kí}\\
 \textsc{th} & \textsc{prvb} & \textit{ń}\\
 \textsc{th} & \textsc{prvb} & \textit{ná}\\
 \textsc{th} & \textsc{prvb} & \textit{ni} (\textit{n})\\
 \textsc{th} & \textsc{prvb} & \textit{ní} (\textit{oo})\\
 \textsc{th} & \textsc{prvb} & \textit{ntsí} (\textit{tsí}, \textit{ntsé}, \textit{tsé}, \textit{ntsá})\\
 \textsc{th} & \textsc{prvb} & \textit{tá}\\
 \textsc{th} & \textsc{prvb} & \textit{tsí}\\
 \textsc{th} & \textsc{prvb} & \textit{ya}\\
 \textsc{th} & \textsc{prvb} & \textit{yá}\\
 \textsc{th} & \textsc{prvb} & \textit{yiní} (\textit{ini})\\
 \textsc{th} & \textsc{prvb} & \textit{’á} (\textit{’}, \textit{’í})\\
 \textsc{th} & \textsc{prvb} & \textit{’a} (\textit{’}, \textit{’i}, \textit{’aa}, \textit{’o}, \textit{’ii})\\
 throat & \textsc{prvb} & \textit{zá}\\
 to\_and\_fro & \textsc{prvb} & \textit{naaná} (\textit{naanáá})\\
 \textsc{trnsl} & \textsc{qual} & \textit{yii} (\textit{i}, \textit{ii}, \textit{oo})\\
 two &\textsc{prvb} & \textit{n}\\
 up, upward & \textsc{prvb} & \textit{ná} (\textit{ń}, \textit{ní}, \textit{né})\\
 useful & \textsc{prvb} & \textit{cho} (\textit{ch})\\
\lspbottomrule
\end{xltabular}